%%%%%%%%%%%%%%%%%%%%%%%%%%%%%%%%%%%%%%%%%%%%%%%%%%%%%%%%%%%%%%%%%%
\section{Recommended verification suite}\label{sec:recommended_suite}
%%%%%%%%%%%%%%%%%%%%%%%%%%%%%%%%%%%%%%%%%%%%%%%%%%%%%%%%%%%%%%%%%%

\todo{WORKING PROPOSAL. Final membership is not frozen; decisions marked below depend on outstanding evidence (paper-data/manuscript\_todos.md).}

Sections~\ref{sec:benchmark_suite} and~\ref{sec:cross_benchmark} suggest that no single benchmark verifies a continuum FSI code, and that the most widely used benchmarks are not the most discriminating.
We therefore propose a tiered suite.
The core suite is inexpensive enough to be run routinely, for example before each release or after a change to a discretisation or coupling algorithm, and is built around problems with exact or independently converged references.
The extended suite contains three-dimensional, geometrically complex or expensive problems that test behaviour under more realistic conditions and are run less often.

%%%%%%%%%%%%%%%%%%%%%%%%%%%%%%%%%%%%%%%%%%%%%%%%%%%%%%%%%%%%%%%%%%
\subsection{Core suite}\label{sec:core_suite}
%%%%%%%%%%%%%%%%%%%%%%%%%%%%%%%%%%%%%%%%%%%%%%%%%%%%%%%%%%%%%%%%%%

\begin{table}[htbp]
	\centering
	\footnotesize
	\setlength{\tabcolsep}{3pt}
	\caption{Proposed core suite (working proposal).}
	\label{tab:core_suite}
	\begin{tabular}{p{2.8cm}p{1.0cm}p{6.4cm}p{3.4cm}}
		\toprule
		Case & Ref. & What it verifies & Status \\
		\midrule
		ringAddedMass & A & added-mass coupling over $\mu = 0.1$--$10$; temporal order; solid-independent check via wet/dry ratio & core \\
		womersleyTube & A & viscous wave propagation; amplitude and phase; temporal order & core; open time-order finding \\
		collapsibleChannel & B & large-deformation bending wall; strong added mass; solid discretisation; partitioned robustness & core (serial; $\approx2.5$\,h full sweep) \\
		Hron--Turek FSI1 & B & steady large-deformation coupling & core; third level desirable \\
		mokLidDrivenCavity & C-weak & self-convergence of a tension-dominated membrane; definition sensitivity & core candidate; see note \\
		\bottomrule
	\end{tabular}
\end{table}

The core suite (Table~\ref{tab:core_suite}) covers inviscid and viscous added mass, small and large deformation, bending- and tension-dominated walls, transient, periodic and steady regimes, and all three coupling algorithms, with a total cost of a few hours on a workstation.
The two category-A problems are the backbone, because only they measure the total error and can expose a deficit in formal order.
The category-B problems extend coverage to large deformation, within the reference uncertainty of $\approx0.3\%$ (collapsibleChannel) and $\approx0.2\%$ (FSI1).
For each core case we recommend reporting at least three refinement levels in space and in time, the observed orders of at least one displacement-type and one force- or phase-type QoI, a coupling-algorithm comparison and a coupling-tolerance sweep.

\todo{Decide whether mokLidDrivenCavity belongs in the core suite (cheap; tests a different solid regime; but C-weak and only self-convergence evidence) or in the extended suite. An independent converged COMSOL solution of the group-B definition would justify core membership. Decide also whether blobInTreacle, if promoted, enters the core suite as a temporal-order check.}

%%%%%%%%%%%%%%%%%%%%%%%%%%%%%%%%%%%%%%%%%%%%%%%%%%%%%%%%%%%%%%%%%%
\subsection{Extended and high-cost suite}\label{sec:extended_suite}
%%%%%%%%%%%%%%%%%%%%%%%%%%%%%%%%%%%%%%%%%%%%%%%%%%%%%%%%%%%%%%%%%%

\begin{table}[htbp]
	\centering
	\footnotesize
	\setlength{\tabcolsep}{3pt}
	\caption{Proposed extended suite (working proposal).}
	\label{tab:extended_suite}
	\begin{tabular}{p{2.8cm}p{1.2cm}p{6.2cm}p{3.4cm}}
		\toprule
		Case & Ref. & What it tests & Status / condition \\
		\midrule
		Hron--Turek FSI2, FSI3 & C & self-excited oscillation; lift-amplitude convergence; tolerance sensitivity; mesh motion (FSI2) & needs FSI3 4x \\
		3dTube & C-weak & 3-D wave propagation with added mass; time-discretisation sensitivity & needs level 3 \\
		beamInCrossFlow & C/C-weak & 3-D steady external flow; near-body resolution; 3-D large deformation & provisional; graded mesh + independent reference \\
		Hessenthaler & -- & realistic 3-D geometry and large deformation; robustness and scalability & in development \\
		cavityFlexibleBottom & C-weak & steady internal flow, large deformation & provisional; may be dropped \\
		blobInTreacle & C-weak & temporal order in viscous large-deformation flow & provisional; possibly supplementary \\
		\bottomrule
	\end{tabular}
\end{table}

The extended suite (Table~\ref{tab:extended_suite}) is not intended to establish formal order, which the core suite does more cheaply and more reliably, but to test whether the accuracy demonstrated there carries over to three dimensions, to geometrically complex domains, to large fluid-mesh motion and to QoIs (oscillatory forces, near-wall-dominated forces) that converge slowly.
The Hron--Turek problems belong here despite being two-dimensional, because their slowest QoIs require meshes and run times well beyond those of the core suite, and because their published references carry approximately $1\%$ uncertainty.
For these cases we recommend that published values be quoted together with their source (e.g.\ Featflow table or summary value) and reference category, and that agreement with them not be presented as verification without a self-convergence study.

\todo{Final extended-suite membership depends on: FSI3 4x; 3dTube level 3; beamInCrossFlow graded-mesh study and independent reference; Hessenthaler study; cavityFlexibleBottom resolution. A case without at least three mesh levels should not be listed as ``verified'' in any final suite table.}

%%%%%%%%%%%%%%%%%%%%%%%%%%%%%%%%%%%%%%%%%%%%%%%%%%%%%%%%%%%%%%%%%%
\subsection{Executable benchmark infrastructure}\label{sec:infrastructure}
%%%%%%%%%%%%%%%%%%%%%%%%%%%%%%%%%%%%%%%%%%%%%%%%%%%%%%%%%%%%%%%%%%

A benchmark is useful for verification only if its definition, its reference and the procedure used to compare against the reference are reproducible.
In the present work each case is distributed with an executable verification driver, and we recommend the following properties for any FSI verification suite.

\begin{itemize}
	\item \emph{Separation from regression testing.} Regression tests check that results have not changed; verification studies check convergence towards a reference. They have different costs, tolerances and purposes, and are kept as separate, opt-in drivers (\texttt{Allverify}) that never modify the tutorial case itself.
	\item \emph{Complete studies, not single runs.} Each driver runs the mesh, time-step and coupling studies, computes the QoIs and observed orders, and exits with a non-zero status unless every acceptance criterion is met. Acceptance tolerances are set from the measured convergence and the model assumptions, and are stored with the reference rather than in the driver.
	\item \emph{Machine-readable outputs.} Per-run QoIs are written to CSV, with a human-readable summary; diverged, truncated or non-periodic runs are reported as failures rather than read as results.
	\item \emph{Reference provenance.} Reference values are stored in machine-readable files that record, for each value, the publication, figure or table, the extraction method (tabulated, vector-extracted, raster-digitised) and its estimated uncertainty; where possible, the script or source code that generates the reference (collocation solutions, oomph-lib driver) is stored with it.
	\item \emph{Versioning and fingerprints.} Each run records the case-file hashes, the OpenFOAM version and a fingerprint of the solver executable and libraries, and previous runs are reused only if all of these match.
	\item \emph{Archival.} The case set-ups, drivers, raw outputs and reference files used for a publication are archived with a persistent identifier.
\end{itemize}

Such infrastructure allows a verification claim to be re-executed and checked by others rather than accepted on trust, and allows the suite to be re-run as the code changes.
Machine-executable verification with unambiguous pass/fail criteria is also increasingly useful for automated scientific-software development workflows, in addition to conventional human-led development, because it provides an objective check that does not depend on visual comparison of plots.
\todo{Not all drivers currently implement every property above (e.g.\ executable/library fingerprints are recorded by some drivers only; per-level CSV is not yet stored for beamInCrossFlow). Audit before submission and state precisely which properties hold for each case. Add Zenodo DOI.}
